\section{Introdu\c{c}\~ao}
\label{sec:introducao}

As audiências públicas das comissões da Câmara dos Deputados reúnem parlamentares, representantes
do governo, especialistas e entidades da sociedade civil em torno de propostas legislativas e
políticas públicas. Suas transcrições registram o que cada participante defendeu e têm em média 18
mil palavras no PublicHearingBR~\citep{fernandes2024publichearingbr}. O público costuma conhecer as
audiências por matérias jornalísticas que têm em média 4,2\% do tamanho da transcrição e citam, em
média, 10,7 opiniões de 5,2 participantes. O conjunto de dados registra, num resumo estruturado, o
que a matéria atribui a cada pessoa, mas não o ponto da fala de onde cada opinião foi extraída.

Este trabalho investiga se um perfil escrito somente com as falas anteriores de uma pessoa ajuda um
grande modelo de linguagem (LLM) a reconhecer o que a imprensa lhe atribuiu numa audiência
posterior. Responder a essa pergunta exige localizar a fala que sustenta cada opinião publicada,
representar, a partir da fala, o que a pessoa defende entre audiências e declarar a base de cada
estimativa de como ela se posicionaria. Uma estimativa desse tipo pode apoiar a preparação de
audiências e ferramentas de transparência que informem o que um participante já defendeu sobre um
tema.

LLMs usados para simular pessoas tendem, porém, a produzir respostas estereotipadas e a apagar a
diversidade interna dos grupos~\citep{bisbee2024synthetic,wang2025large}. Por isso, a simulação
proposta recebe como material sobre a pessoa, além do nome e do cargo, apenas as falas dela nas
audiências de treino, não fornece ao modelo o partido registrado e declara a base de cada estimativa.
Diferentemente dos trabalhos que condicionam LLMs a perfis sociodemográficos, a opiniões passadas ou
a entrevistas de respondentes de questionário~\citep{argyle2023outofone,hwang2023aligning,park2026llm},
ela trata de pessoas públicas identificadas e é avaliada com o que a imprensa lhes atribuiu em
audiências posteriores às usadas no perfil.

O primeiro dos três componentes da abordagem (Seção~\ref{sec:metodo}) é a Unidade Deliberativa
Verificável (UDV), que procura na fala da própria pessoa uma sentença que sirva de evidência
candidata para cada opinião do resumo estruturado. A UDV registra os offsets de caracteres da
sentença escolhida, um nível com o tipo de vínculo e a proveniência, que indica a origem do vínculo.
Os perfis por ator, segundo componente, são escritos por um LLM com as falas de cada pessoa nas
audiências de treino, em blocos de posições, critérios e valores, alinhamentos declarados e forma de
argumentar. O terceiro, a simulação, fornece o perfil a um LLM e acrescenta trechos das falas
anteriores da pessoa recuperados pelo assunto. Nas perguntas de múltipla escolha, e não na geração
aberta, a simulação combina as probabilidades das quatro opções com e sem esse material por
\textit{classifier-free guidance} (CFG), o que desfavorece as opções que o modelo já escolheria só
com o nome e o cargo.

As UDVs das audiências de validação e de teste ligam as opiniões publicadas aos atores pelo turno da
evidência, e cada opinião ligada a um ator com perfil gera uma pergunta de múltipla escolha. A opção
correta é a opinião na redação do resumo estruturado, e a sentença de evidência não aparece na
pergunta. As perguntas de validação escolhem dois hiperparâmetros da simulação, e as de teste medem o
efeito do perfil, dos trechos e da CFG. O trabalho se inspira nas Trilhas A e B do desafio: a UDV e a
recuperação restrita à fala da pessoa tratam da auditoria de documentos longos, e os perfis e a
simulação, da mineração de opiniões e de atores.

\medskip\noindent\textbf{Contribui\c{c}\~oes.} \textbf{(i)} A UDV localiza na fala da própria pessoa
evidência candidata para 2.105 das 2.203 opiniões do PublicHearingBR, declara o nível e a
proveniência de cada vínculo e tem a consistência dos registros conferida por um verificador
separado. \textbf{(ii)} Os perfis por ator são escritos por um LLM com instruções que proíbem
atribuir à pessoa o que vem do assunto da audiência, de terceiros, do partido ou de conhecimento
externo. \textbf{(iii)} O método de simulação de atores declara a base de cada estimativa e
acompanha a fala simulada de uma justificação cujas referências ao material são conferidas
automaticamente. \textbf{(iv)} O protocolo de avaliação usa como gabarito opiniões publicadas sobre
audiências posteriores às do perfil, ligadas ao ator pelo turno da evidência da UDV, com distratores
da mesma audiência. O protocolo mede separadamente o efeito do perfil, dos trechos e da CFG; nele,
<resultado principal>.
